\documentclass[11pt]{article}
\usepackage[margin=27mm]{geometry}
\usepackage{amsmath,amssymb,amsthm,booktabs,hyperref}
\newtheorem{theorem}{Theorem}
\newtheorem{lemma}[theorem]{Lemma}
\newtheorem{corollary}[theorem]{Corollary}
\newcommand{\PSL}{\operatorname{PSL}}
\newcommand{\PGL}{\operatorname{PGL}}
\newcommand{\N}{\operatorname{N}}
\newcommand{\F}{\mathbb F}
\title{Two-transitive coset actions in dimension two}
\author{Human Compression Project}\date{}
\begin{document}\maketitle
\begin{abstract}
We derive the two-transitive actions needed in the two fixed Part II
originals from the classical subgroup classification of the linear
groups of dimension two. A point stabilizer of order $b$ has degree
at most $b+1$. This bound, together with a short divisibility argument,
removes the subfield alternatives and leaves a small list of exceptional
actions. The resulting common argument also determines the possible
overgroups of the small intersections used in those originals. We do
not assume their Proposition 4.6 or its unpublished reference, and
we do not count embeddings or pairs of parabolics here.
\end{abstract}

\section{Premises and the elementary bound}
We write $D_{2n}$ for a dihedral group of order $2n$. The classical
subgroup theorems of Dickson, Moore and Suzuki are borrowed premises:
the proper subgroups of $\PSL(2,t)$ are cyclic or dihedral groups,
elementary abelian groups and their scalar cyclic extensions, $A_4$,
$S_4$, $A_5$, and the subfield groups $\PSL(2,s)$ and $\PGL(2,s)$.
We use their usual field restrictions, maximal-subgroup lists and
normalizers. In particular a subgroup $\PGL(2,s)$ in $\PSL(2,t)$,
for odd $s$, requires an even extension degree. The normalizer of a
four-group in odd characteristic contains an $A_4$. A torus has its
dihedral normalizer. The two $S_4$ classes in $\PSL(2,7)$ and the
two $A_5$ classes in $\PSL(2,9)$ and $\PSL(2,11)$ are interchanged
by $\PGL(2,t)$. These are subgroup and fusion premises, not assertions
about the geometries to be recovered. They are identified separately
in Sections 4.1, 4.2 and 4.4 of the fixed published source
\cite{source}. We retain that source's attribution to
\cite{dickson,moore,suzuki}; we do not represent the historical
books as newly checked source editions. The modern statement in
\cite[Theorem 2 and Lemmas 1--2]{cameron}, consulted separately,
also records the subgroup types in odd characteristic and the
small exceptional subgroup lattices. Its entire design theory is
not an input to this proof.

We also use the classical isomorphisms
$\PSL(2,4)\cong\PSL(2,5)\cong A_5$,
$\PSL(2,7)\cong\operatorname{GL}(3,2)$,
$\PSL(2,9)\cong A_6$ and $\PGL(2,5)\cong S_5$.
The simplicity of $\PSL(2,t)$ for $t\geq4$ is retained as a premise.
The small subgroup lattices of $A_4,S_4,A_5$ also follow directly
from their natural permutation actions; we use their maximal
subgroup lists below.

\begin{lemma}\label{bound}
Suppose that $H$ acts two-transitively on the cosets of a proper
subgroup $B$, and put $d=[H:B]$. Then $B$ is maximal in $H$,
$d-1$ divides $|B|$, and $|H|\leq |B|(|B|+1)$.
\end{lemma}
\begin{proof}
A nontrivial system of blocks would prevent a point stabilizer from
being transitive on the other points. Thus the action is primitive,
and $B$ is maximal. The $B$-orbit on the other $d-1$ points has
length $d-1$, which divides $|B|$. Multiplying $d\leq |B|+1$
by $|B|$ proves the bound. Faithfulness is unnecessary.
\end{proof}

\section{The simple groups}
\begin{theorem}\label{simple}
Let $H=\PSL(2,t)$, where $t\geq4$ is a prime power. The
two-transitive coset actions are the natural action with point
stabilizer $E_t:C_{(t-1)/(2,t-1)}$, and the following exceptional
actions, with repetitions under the small isomorphisms suppressed:
\[
\begin{array}{c|c|c}
H&B&[H:B]\\\hline
A_5&A_4&5\\
A_5&D_{10}&6\\
\PSL(2,7)&S_4&7\\
\PSL(2,9)&A_5&6\\
\PSL(2,11)&A_5&11.
\end{array}
\]
Here each embedding class of the displayed stabilizer gives its
corresponding coset action. The assertion is not a count of ambient
embeddings into a larger group.
\end{theorem}
\begin{proof}
The kernel is normal in the simple group $H$ and is contained in
the proper subgroup $B$, hence is trivial. Apply the maximal-subgroup
list and Lemma~\ref{bound}.

For a torus normalizer the two possible degrees are
$t(t-1)/2$ and $t(t+1)/2$. For odd $t\geq7$, their point
stabilizers have order at most $t+1$, whereas
$t(t-1)/2-1>t+1$. For even $t\geq8$, their order is at most
$2(t+1)$ and $t(t-1)/2-1>2(t+1)$. Thus only $t=4,5$
remain. In $A_5$, $D_{10}$ has index six; $D_6$ has index ten,
which fails $9\mid6$.

If $B=A_4$, the order bound is $|H|\leq156$, so $t\leq6$
in odd characteristic and $t\leq5$ in even characteristic.
Only $t=4,5$ give the required subgroup. If $B=S_4$, the bound
is $|H|\leq600$. The only possible fields with an $S_4$
subgroup are $t=7,9$; for $t=9$ the degree is fifteen and
$14\nmid24$. If $B=A_5$, the bound is $|H|\leq3660$.
Apart from $H=A_5$, the possible orders divisible by sixty are
those for $t=9,11,19$; the last has degree fifty-seven and
$56\nmid60$. These finite checks use the order
$t(t^2-1)/(2,t-1)$ and the prime powers in the stated bounds.

For an even field, a maximal proper subfield group has $t=s^r$
with $r$ prime. At $r=2$ its index is $s(s^2+1)$ and its
order is $s(s^2-1)$, so the index minus one already exceeds
its order. The index increases with $r$. For an odd field and
a maximal simple subfield group, $r\geq3$. Its index is at least
$s^2(s^4+s^2+1)$, again exceeding its order plus one. The
remaining maximal subfield possibility is $\PGL(2,s)$ in
$\PSL(2,s^2)$ for odd $s\geq3$. Its order is
$b=s(s^2-1)$ and its index is $d=s(s^2+1)/2$. But
\[
 1<\frac{b}{d-1}
   =2-\frac{4s-4}{s^3+s-2}<2.
\]
Consequently $d-1$ cannot divide $b$. This removes every proper
subfield alternative.

The listed actions do exist. Translations in a Borel subgroup are
transitive on the finite points of the projective line, proving the
natural case. The natural $A_5$ action on five letters is
two-transitive. Its action on the six Sylow five-subgroups is
two-transitive because a Sylow five-subgroup fixes only itself:
another fixed point would put it in a different normalizer of
order ten. It therefore has one orbit on the other five points.
The action of $\operatorname{GL}(3,2)$ on its seven nonzero vectors
is two-transitive, since two distinct such vectors are independent.
The natural $A_6$ action gives the degree-six case.

For $H=\PSL(2,11)$ choose one class of $A_5$ subgroups. It has
eleven members. There are fifty-five $A_4$ subgroups, all conjugate
and self-normalizing; each chosen $A_5$ contains five. Double
counting shows that each $A_4$ lies in exactly one member of the
chosen $A_5$ class. Hence distinct members do not intersect in
an $A_4$. In the action on that class the stabilizer $A_5$ has no
additional fixed point, because the normalizer of an $A_5$ is itself.
Its possible nontrivial orbit lengths at most ten are five, six
and ten, by the subgroup orders of $A_5$. Length five would give
an $A_4$ intersection and is impossible. Ten remaining points
cannot be partitioned into orbits of length six. They form one
orbit of length ten. This proves two-transitivity, for either
of the two $A_5$ classes.
\end{proof}

\section{The projective groups}
\begin{theorem}\label{projective}
Let $H=\PGL(2,t)$ for odd $t\geq5$. Its two-transitive
coset actions have point stabilizer $\PSL(2,t)$ and degree two,
the natural Borel $E_t:C_{t-1}$ and degree $t+1$, or, when
$t=5$, $S_4$ and degree five. For $\PGL(2,3)=S_4$ the
stabilizers are $A_4,D_8,S_3$, of degrees two, three and four.
\end{theorem}
\begin{proof}
Put $N=\PSL(2,t)$. If $N$ is in the kernel, the nontrivial
action factors through $H/N=C_2$, giving the first case.
Otherwise a nontrivial kernel would meet the simple group $N$
trivially and centralize it. The centralizer of $N$ in $H$ is
trivial: commuting with all translations fixes their unique common
fixed point and gives an affine map with scalar one; commuting
also with the opposite unipotent subgroup makes its translation
zero. Thus the action is faithful.

Primitivity now makes $N$ transitive. For a point stabilizer $B$
write $C=B\cap N$. Then $H=NB$, $|B|=2|C|$, and
$d=[H:B]=[N:C]$. In particular
\begin{equation}\label{twice}
 d-1\mid 2|C|,\qquad |N|\leq |C|(2|C|+1).
\end{equation}
We apply the full subgroup list to $C$, which need not be maximal
in $N$.

If $C$ has nontrivial unipotent radical $O_p(C)$, this radical
is characteristic in $C$ and has a unique common projective fixed
point. Since $C\lhd B$, $B$ fixes this point. Maximality of $B$
in $H$ makes $B$ the full Borel, giving the natural action.
This covers all affine types, including nontrivial unipotent
cyclic subgroups.

If $C$ is a nontrivial semisimple cyclic group, $B$ normalizes its
torus. For an involution this follows by centralizing it, which
preserves its two eigenlines as an unordered pair. If $C$ is
dihedral with rotation group of order greater than two, that
rotation group is characteristic, with the same conclusion.
Thus maximality forces $B$ to be a torus normalizer. Its order
is at most $2(t+1)$ and the smaller possible degree is
$t(t-1)/2$. For $t\geq7$, $d-1>2(t+1)$. At $t=5$ the
two degrees are ten and fifteen, with stabilizer orders twelve
and eight, and both fail divisibility. The case $C=1$ fails
\eqref{twice}. If $C$ is a four-group, its normalizer contains
an $A_4$, while $|B|=8$. This normalizer is proper, since
$C$ cannot be normal in the simple group $N$. It would strictly
contain the maximal subgroup $B$, a contradiction.

If $C=A_4$, \eqref{twice} gives $|N|\leq300$. Only
$t=5,7$ can occur; $t=7$ gives $d=14$ and $13\nmid24$.
For $t=5$ the extension is $S_4$ inside $S_5$, with its natural
degree-five action. If $C=S_4$, the bound is $|N|\leq1176$.
The possible fields are $t=7,9$; $t=9$ fails $14\mid48$.
At $t=7$, an element of $B\setminus N$ would normalize the
$S_4$ subgroup, contrary to the interchange of its two $N$-classes
by $H$. If $C=A_5$, the bound is $|N|\leq7260$.
The proper possibilities are $t=9,11,19$; $t=19$ fails
$56\mid120$. For $t=9,11$, interchange of the two $A_5$
classes again excludes such a $B$. The case $C=N=A_5$
belongs to the already separated degree-two action.

Finally suppose $C$ is a proper subfield group with field size
$s$. For $C=\PSL(2,s)$ the smallest possible index occurs
at $t=s^2$ and equals $s(s^2+1)$; this already exceeds
$2|C|+1=s(s^2-1)+1$. For $C=\PGL(2,s)$, $t$ is an
even power of $s$. At $t=s^2$,
\[
 3<\frac{2|C|}{d-1}
   =4-\frac{8s-8}{s^3+s-2}<4 \qquad(s\geq3),
\]
so divisibility fails. At any larger allowed extension degree,
even the index for $t=s^3$ is
$s^2(s^4+s^2+1)/2>2s(s^2-1)+1$, and the bound fails.
This exhausts $C$.

The Borel and degree-two actions plainly exist, and the exceptional
degree-five action is the natural $S_5$ action. The maximal
subgroups $A_4,D_8,S_3$ of $S_4$ give respectively its quotient
action on two points, its $S_3$ quotient on three points, and its
natural action on four points. They exhaust the case $t=3$.
\end{proof}

\section{Solvable overgroups}
The following elementary observations handle the remaining Dickson
types without a classification of general two-transitive groups.

\begin{lemma}\label{solvable}
For a dihedral group, a two-transitive coset action has degree two
or three. In degree two the stabilizer is any subgroup of index
two; in degree three it is a dihedral subgroup of index three.

Let $H=V:C$ be a scalar affine subgroup in characteristic $p$,
where $V\ne0$ is elementary abelian and $C$ is cyclic of order
prime to $p$. A two-transitive coset action either has degree
two and its stabilizer contains $V$, or has a stabilizer conjugate
to $W:C$, with $W$ a $C$-invariant subgroup of $V$ and
\[
 [V:W]=|C|+1.
\]
The induced scalar action on $V/W\setminus\{0\}$ is regular.
\end{lemma}
\begin{proof}
For a dihedral group, the orbits of its normal rotation group
form blocks. If this group is in the kernel, the action factors
through $C_2$. Otherwise it is transitive in the faithful quotient.
A transitive abelian permutation group is regular, since all
point stabilizers coincide and hence lie in the kernel. The
point stabilizer in the resulting dihedral permutation group
therefore has order at most two. Transitivity on the other
points implies that the degree is at most three. A subgroup of
index three with a transitive rotation group contains a reflection
and is dihedral. Conversely the quotients on two and three points
are $C_2$ and $S_3$, respectively.

For the affine group, primitivity likewise says that $V$ is in
the kernel or is transitive in the faithful quotient. In the
first case the action factors through the abelian group $C$;
two-transitivity then forces degree two. In the second case put
$W=V\cap B$. Regularity of the image of $V$ shows that $W$
is the kernel within $V$, hence is $H$-invariant, and $H=VB$.
The image of $B$ in $H/V=C$ is all of $C$. Modulo $W$,
the subgroup $B/W$ is a complement to $V/W$. Such complements
are conjugate by $V/W$: one may average an affine fixed point
over the cyclic complement, since $|C|$ is invertible in
characteristic $p$. Thus $B$ is conjugate to $W:C$.

For a nonidentity scalar $c$, multiplication by $c-1$ is
invertible on $V$ and maps the invariant finite subspace $W$
onto itself. It is therefore invertible on $V/W$. Thus the
nonidentity elements of $C$ have no nonzero fixed vector in
$V/W$, and two-transitivity says that this free action on the
nonzero vectors is also transitive. This proves the formula.
\end{proof}

\section{The small intersections in the fixed originals}
\begin{corollary}\label{small}
Within the field scopes used in the published Part II source,
the possible proper two-transitive overgroups of the indicated
intersection are as follows. Each listed dihedral overgroup
requires its usual torus divisibility condition.
\[
\begin{array}{c|l|l}
B&\text{ambient field scope}&H\\\hline
D_{10}&\text{the maximal }A_5\text{ cases}&D_{20},D_{30},A_5\\
A_4&\text{the maximal }A_5\text{ cases}&S_4,A_5,E_{16}:C_3\\
D_6&q\text{ prime},\ q\equiv\pm1\pmod8&D_{12},D_{18},S_4\\
D_8&q\text{ prime},\ q\equiv\pm1\pmod8&D_{16},D_{24},S_4\\
A_4&q\text{ prime},\ q\equiv\pm1\pmod8&S_4,A_5\\
A_5&q=9^m\text{ or }11^m,\ m\text{ odd prime}
 &\PSL(2,9)\text{ or }\PSL(2,11)\\
S_4&q=7^m,\ m\text{ odd prime}&\PSL(2,7)\\
D_8&q=9&S_4\\
S_4&q=25&\PGL(2,5).
\end{array}
\]
The last two alternatives in the $A_5$ row are restricted to
their actual defining characteristic. The affine $E_{16}:C_3$
alternative occurs only in characteristic two and needs at
least two dimensions over $\F_4$ in the ambient additive field.
An $S_4$ or $A_5$ alternative is included only when it embeds
in the specified ambient group.
\end{corollary}
\begin{proof}
For the cyclic and dihedral Dickson types, a nonabelian
stabilizer cannot lie in a cyclic group. Lemma~\ref{solvable}
gives exactly the displayed double and triple dihedral orders.
For $q=9$ neither $D_{16}$ nor $D_{24}$ embeds: the maximal
tori have orders four and five.

Consider an affine overgroup. Its stabilizer either contains
the whole additive kernel with index two, or is $W:C$ as in
Lemma~\ref{solvable}. For $B=A_4$, the latter forces
$W=E_4$, $C=C_3$, and $[V:W]=4$, giving $E_{16}:C_3$.
An index-two extension would instead require a faithful scalar
$C_6$ on $E_4$, impossible since its scalar group has order
three. This also shows directly that no affine alternative
exists in odd characteristic for $B=A_4$.

For $B=D_{10}$, a nontrivial additive normal subgroup must be
$C_5$. The transitive-kernel alternative would demand an additive
quotient of size $|C_2|+1=3$ in characteristic five, impossible.
The index-two alternative would be $C_5:C_4$. In the present
maximal-$A_5$ scope a characteristic-five field has size $5^m$
with $m$ odd, so its scalar torus has order $(5^m-1)/2$, which
is twice an odd number. It contains no $C_4$. For $B=D_6$
the affine possibility would require characteristic three;
the stipulated prime $q\equiv\pm1\pmod8$ is not three.
For $D_8$ in odd characteristic there is no nontrivial normal
subgroup of odd prime-power order. Finally $S_4$ and $A_5$
cannot occur as these affine stabilizers: in odd characteristic
they have no nontrivial normal subgroup of defining prime-power
order, and they are not cyclic. Thus no affine cases were lost.

The small exceptional overgroups are read directly from the
maximal-subgroup lists: $A_4$ has its degree-four action with
stabilizer $C_3$; $S_4$ has the three actions in
Theorem~\ref{projective}; $A_5$ has the degree-five and degree-six
actions in Theorem~\ref{simple}. These give precisely the
$S_4$ and $A_5$ entries above.

For all larger subfield overgroups, apply Theorems~\ref{simple}
and~\ref{projective} to the abstract overgroup. Reading off
the displayed stabilizers gives $\PSL(2,7)$ above $S_4$,
$\PSL(2,9)$ and $\PSL(2,11)$ above $A_5$, and
$\PGL(2,5)$ above $S_4$ or its normal $A_5$. Such groups
are neither cyclic, dihedral, affine nor one of the three small
exceptional Dickson types, so their embeddings must be subfield
embeddings in the actual defining characteristic. This excludes
$\PGL(2,5)$ from the characteristic-three and -eleven $A_5$
rows, and from the characteristic-seven $S_4$ row. It excludes
$\PSL(2,7)$ in characteristic five. The natural Borel cases
give no further small stabilizers. This proves the table.
\end{proof}

\section{Source use and limits}
Corollary~\ref{small} proves the overgroup possibilities used
in Lemmas 5.1, 5.2, 5.7--5.9, 5.14, 5.15, 5.17 and 5.20
of the fixed published original L063. It supplies the common
argument for Lemmas 12, 17--19, 24, 25 and 27 of the retained
appendix L064, with the published parameter restrictions stated
explicitly. The retained appendix begins at Lemma 12; it does
not contain the advertised proof of published Lemma 5.1.

Theorem~\ref{simple} and Theorem~\ref{projective} replace the
simple and projective portions of the unpublished classification
cited as Proposition 4.6. They may be used in the separate
generic subfield argument. The cyclic, dihedral and affine
arguments explain their solvable counterparts, but no claim is
made here that every literal assertion about uniqueness or class
fusion in that proposition is correct. In particular a
representation of an abstract subgroup and an embedding of that
subgroup in a larger group are different counting problems.

The small bounds are finite integer calculations, not a group
census. The classical subgroup classifications, their field
restrictions, the stated exceptional isomorphisms and fusion facts
remain attributed prior results. This manuscript supplies the
common elimination and the elementary two-transitivity arguments.
It does not establish all configuration counts, complete either
original, solve an open conjecture, or claim Lean verification
or a compression factor.
\begin{thebibliography}{9}
\bibitem{source} F.~Buekenhout, J.~De Saedeleer and D.~Leemans,
On the rank two geometries of the groups $\PSL(2,q)$: part II,
\emph{Ars Mathematica Contemporanea} \textbf{6} (2013), 365--388.
The separate retained appendix has 21 physical pages and ends at
Proposition 30.
\bibitem{dickson} L.~E.~Dickson, \emph{Linear Groups: With an
Exposition of the Galois Field Theory}, 1901; later Dover reprint.
\bibitem{moore} E.~H.~Moore, The subgroups of the generalized
finite modular group, \emph{Decennial Publications of the
University of Chicago} \textbf{9} (1904), 141--190.
\bibitem{suzuki} M.~Suzuki, \emph{Group Theory I}, Grundlehren
der mathematischen Wissenschaften 247, Springer, 1982, p.~417.
\bibitem{cameron} P.~J.~Cameron, G.~R.~Omidi and B.~Tayfeh-Rezaie,
$3$-Designs from $\PGL(2,q)$, \emph{Electronic Journal of
Combinatorics} \textbf{13} (2006), R50,
\href{https://doi.org/10.37236/1076}{doi:10.37236/1076}.
\end{thebibliography}
\end{document}
